% !TEX program = lualatex
% Generated by new_homework.sh — Math 104 Homework 3
\documentclass[11pt]{article}

\usepackage{fontspec}
\usepackage{microtype}
\usepackage{mathtools,amssymb,amsfonts,amsthm,bm}
\usepackage{enumitem}
\usepackage{booktabs,array,xcolor}
\usepackage{geometry,fancyhdr,setspace}
\usepackage{pdfpages}
\usepackage[hidelinks]{hyperref}
\usepackage[nameinlink,noabbrev]{cleveref}

% ---------- Page and paragraph rhythm (matches lecture/solutions templates) ----------
\geometry{margin=1in,headheight=15pt}
\setstretch{1.12}
\setlength{\parindent}{0pt}
\setlength{\parskip}{0.55em plus 0.12em minus 0.08em}
\setlength{\abovedisplayskip}{0.8em plus 0.2em minus 0.1em}
\setlength{\belowdisplayskip}{0.8em plus 0.2em minus 0.1em}
\allowdisplaybreaks[3]
\setlist{leftmargin=*,topsep=0.35em,itemsep=0.15em,parsep=0pt}

% ---------- Metadata ----------
\newcommand{\studentname}{Keshav Ramamurthy}
\newcommand{\coursename}{Math 104}
\newcommand{\courselongname}{Real Analysis}
\newcommand{\assignment}{Homework 3}
\newcommand{\term}{Fall 2026}

% ---------- Core notation (same as ~/.config/nvim/textemplate.tex) ----------
\newcommand{\N}{\mathbb{N}}
\newcommand{\Z}{\mathbb{Z}}
\newcommand{\Q}{\mathbb{Q}}
\newcommand{\R}{\mathbb{R}}
\newcommand{\C}{\mathbb{C}}
\newcommand{\F}{\mathbb{F}}
\newcommand{\K}{\mathbb{K}}
\newcommand{\E}{\mathbb{E}}
\newcommand{\Prob}{\mathbb{P}}
\newcommand{\1}{\mathbf{1}}
\newcommand{\eps}{\varepsilon}
\newcommand{\dd}{\,\mathrm{d}}
\newcommand{\abs}[1]{\left\lvert#1\right\rvert}
\newcommand{\norm}[1]{\left\lVert#1\right\rVert}
\newcommand{\ip}[2]{\left\langle#1,#2\right\rangle}
\newcommand{\set}[1]{\left\{#1\right\}}
\newcommand{\ceil}[1]{\left\lceil#1\right\rceil}
\newcommand{\floor}[1]{\left\lfloor#1\right\rfloor}
\newcommand{\given}{\,\middle\vert\,}
\newcommand{\indep}{\perp\!\!\!\perp}
\newcommand{\wh}[1]{\widehat{#1}}
\newcommand{\deriv}[2]{\frac{\mathrm{d}#1}{\mathrm{d}#2}}
\newcommand{\pderiv}[2]{\frac{\partial#1}{\partial#2}}
\newcommand{\lto}{\longrightarrow}
\newcommand{\st}{\text{ such that }}
\DeclareMathOperator{\Aut}{Aut}
\DeclareMathOperator{\End}{End}
\DeclareMathOperator{\Gal}{Gal}
\DeclareMathOperator{\Hom}{Hom}
\DeclareMathOperator{\Ker}{ker}
\DeclareMathOperator{\im}{im}
\DeclareMathOperator{\rank}{rank}
\DeclareMathOperator{\nullity}{nullity}
\DeclareMathOperator{\Span}{span}
\DeclareMathOperator{\tr}{tr}
\DeclareMathOperator{\diag}{diag}
\DeclareMathOperator{\spec}{spec}
\DeclareMathOperator{\ord}{ord}
\DeclareMathOperator{\supp}{supp}
\DeclareMathOperator{\diam}{diam}
\DeclareMathOperator{\Var}{Var}
\DeclareMathOperator{\Cov}{Cov}
\DeclareMathOperator*{\argmin}{arg\,min}
\DeclareMathOperator*{\argmax}{arg\,max}

% ---------- Statement environments ----------
\newtheorem{problem}{Problem}
\newtheorem{theorem}{Theorem}
\newtheorem{lemma}{Lemma}
\newtheorem{proposition}{Proposition}
\newtheorem{corollary}{Corollary}
\theoremstyle{definition}
\newtheorem{definition}{Definition}
\newtheorem{example}{Example}
\newtheorem{remark}{Remark}

% ---------- Solution machinery (matches solutions.tex / textemplate.tex) ----------
\newenvironment{solution}{\begin{proof}[Solution]}{\end{proof}}
\newenvironment{answer}{\begin{proof}[Answer]}{\end{proof}}
\renewcommand{\qedsymbol}{$\square$}
\newcommand{\exercise}[1]{%
  \par\goodbreak\medskip\noindent\textbf{Exercise #1}\par\nopagebreak\smallskip\nopagebreak}
\newcommand{\exref}[1]{\textsf{\textbf{Exercise~#1}}}
\newenvironment{claim}{\par\smallskip\noindent\textit{Claim.}\ }{\par\smallskip}
\newenvironment{idea}{\par\smallskip\noindent\textit{Idea.}\ }{\par\smallskip}
\newenvironment{proofsketch}{\par\smallskip\noindent\textit{Proof sketch.}\ }{\par\smallskip}

% ---------- Header ----------
\pagestyle{fancy}
\fancyhf{}
\fancyhead[L]{\coursename}
\fancyhead[C]{\assignment\ — my solutions}
\fancyhead[R]{\studentname}
\fancyfoot[C]{\thepage}

\begin{document}

% ---------- Title ----------
\begin{center}
  {\Large\sffamily\bfseries \coursename\ -- \courselongname}\par\vspace{3pt}
  {\sffamily \assignment\ — my solutions}\par\vspace{3pt}
  {\small\sffamily \term\quad\textbullet\quad \studentname}
\end{center}
\medskip
\hrule
\bigskip



\section*{Solutions}
\subsubsection*{Problem 1}
\begin{solution}
    We start by dividing off by $n^6$ to hopefully manipulate this nicely:
    $$\frac{11n^6-4n^5+9n^2-7}{13n^6+5n^4+2n+1} = \frac{11 - 4\frac{1}{n} +
        9\frac{1}{n^4} - 7\frac{1}{n^6}}{13 + 5\frac{1}{n^2} + 2\frac{1}{n^5} +
    \frac{1}{n^6}}$$
    Let's examine the numerator: 

    Notice, $\lim_{n\rightarrow\infty} \left(11 - 4\frac{1}{n} +
        9\frac{1}{n^4} - 7\frac{1}{n^6} \right)= 11$ because 
        $$\lim_{n\rightarrow\infty}\left(11 - 4\frac{1}{n} + 9\frac{1}{n^4} -
        7\frac{1}{n^6} \right) = \lim 11 -4 \lim\frac{1}{n} + 9\lim \frac{1}{n^4} - 7
        \lim \frac{1}{n^6} = 11  $$
    using $\lim_{n\rightarrow \infty}
    \frac{1}{n^p} = 0$ for $p > 0$.
    This exact same argument can be applied for the denominator where everything but the
    13 goes to zero because:
    $$\lim_{n\rightarrow\infty} \left(13 + 5\frac{1}{n^2} + 2 \frac{1}{n^5} + \frac{1}{n^6}\right) =
    \lim 13 + 5\lim \frac{1}{n^2} + 2\lim\frac{1}{n^5} + \lim \frac{1}{n^6} = 13 +
    5(0) + 2(0) + 0 = 13$$
    And the quotient rule applies because the denominator limit is nonzero.
    Thus, this limit evaluates to $\frac{11}{13}$.
\end{solution}

\subsubsection*{Problem 2}
\begin{solution}
We are given the following: $$\lim_{n\rightarrow\infty} x_n = -2, \lim_{n\rightarrow\infty}
y_n = 5$$ and that $x_n + 2y_n \ne 0$ for every $n \in \mathbb{N}$.
\par Because the denominator will always be nonzero as given, we can consider separately the
limit of the numerator and denominator, and apply the quotient rule because the limit of
the denominator is also nonzero.

\par When considering $$\lim_{n\rightarrow\infty}(2y_n-x_n) = 2\lim y_n - \lim x_n = 2(5) - (-2) = 12. $$
To consider the denominator, notice $$\lim_{n\rightarrow \infty} (x_n+2y_n)^2 = \left(\lim_{n\rightarrow \infty} (x_n + 2y_n)\right)^2$$
from the product rule, and so we have that this limit is equal to 
$$\left(\lim x_n + 2\lim y_n\right)^2 = (-2 + 10)^2 = 64$$
So our answer is $\frac{12}{64} = \frac{3}{16}$.
\end{solution}

\subsubsection*{Problem 3}
\begin{solution}
    We first notice that
    $$\lim_{n\rightarrow\infty}(a_n^2+b_n^2+1)=a^2+b^2+1\ge1>0,$$
    so the limit of the denominator is nonzero.
    \par So in order to continue, let's consider the numerator and denominator
    separately.
    We have $$\lim_{n\rightarrow\infty}( a_n^4 - 2a_nb_n + 3) =
    \lim a_n^4 - 2\lim a_n \lim b_n + \lim 3 = a^4 - 2ab + 3$$
    Considering the denominator similarly, notice that 
    $$\lim_{n\rightarrow\infty} (a_n^2 + b_n^2 + 1) = (\lim a_n)^2 + (\lim b_n)^2 + 1 =a^2 + b^2 + 1$$
    so we are done, as we have shown that $$\lim_{n\rightarrow\infty} u_n =
    \frac{a^4-2ab+3}{a^2+b^2+1}.$$
\end{solution}

\subsubsection*{Problem 4}
\begin{solution}
    We have that $x_1 = 1$ and $x_{n + 1} = x_n (4x_n - 1).$
    Notice we can write:
    We claim $x_n \ge 3^{n-1}$ for all positive integers $n$, and will show this by induction.
\par The base case is trivial as $x_1 = 1 \ge 3^{1-1} =1$.
\par For the inductive step, assume that $x_k \ge 3^{k-1}$ for some positive integer $k$.
Then, $$x_{k+1} = x_k (4x_k - 1).$$
Let's look at our inductive hypothesis:
$$x_k \ge 3^{k-1}\ge 1 \rightarrow 4x_k \ge 4 
\rightarrow 4x_k - 1 \ge 3.$$
So, we have that $4x_k - 1 \ge 3$ and we are done because then
$$x_{k + 1} \ge 3^{k - 1}\cdot 3 = 3^k.$$
So by induction this is true for all $k \in \mathbb{Z}^{+}$.
We have that $x_n \ge 3^{n-1}$ and $\lim_{n\rightarrow\infty} 3^{n-1}=+\infty$,
and so by comparison, $\lim_{n\rightarrow \infty}x_n = +\infty$.
\end{solution}

\subsubsection*{Problem 5}
\begin{solution}
    We have that $r_n \rightarrow 0$ and $\exists N_0 \in \mathbb{N}$ such that
    $$s_n \le t_n + r_n \text{ for every } n > N_0.$$
    We write $t_n \ge s_n - r_n$ and just wish to show rigorously that "$s_n$" gets
    large and $r_n$ is very small.
    
    Consider an arbitrary $M > 0$. Since $s_n \rightarrow +\infty$, there exists
    $N_2 \in \mathbb{N}$ such that for every $n > N_2$,
    $$s_n > M + 1.$$
    
    \par We write
    $$t_n \ge s_n - r_n.$$
    
    Along with this, we notice since $r_n \rightarrow 0$, there exists $N_1 \in \mathbb{N}$
    such that $|r_n| < 1$ for every $n > N_1$, so notice that for every
    $n > \max\{N_0,N_1,N_2\}$,
    $$t_n \ge s_n - r_n > s_n - 1 > M + 1 - 1 = M.$$
    
    So for arbitrary $M$ there exists $N = \max\{N_0,N_1,N_2\}$ such that for every $n > N$,
    both $s_n > M + 1$ and $|r_n| < 1$ hold, which satisfies $t_n > M$.
    Therefore,
    $$\lim_{n\to\infty}t_n=+\infty.$$
\end{solution}

\subsubsection*{Problem 6}
\begin{solution}
% Write your proof, calculation, or argument here.
\end{solution}

\subsubsection*{Problem 7}
\begin{solution}
% Write your proof, calculation, or argument here.
\end{solution}

\subsubsection*{Problem 8}
\begin{solution}
% Write your proof, calculation, or argument here.
\end{solution}

\end{document}
